\studytypesubsection{LLMs as Annotators}
\label{sec:llms-as-annotators}

In the annotator role, LLMs perform qualitative coding---the annotation of natural language text such as requirements, interview transcripts, or open-ended survey responses---that researchers would otherwise do by hand.

\studytypeparagraph{Description}
Coding is a time-consuming manual process~\cite{DBLP:journals/ase/BanoHZT24}. LLMs can augment this process, suggest new codes, and label artifacts based on a pre-defined coding guide much faster than humans can~\cite{DBLP:conf/chi/HeHDRH24}. This covers both \emph{closed coding}, where the LLM labels artifacts against a predefined coding guide, and \emph{open coding}, where the LLM proposes new codes from the data. In closed coding, LLM labels can be assessed against those of human coders applying the same guide; in open coding, researchers review the resulting codebook for adequacy (e.g., level of abstraction, redundancies between codes).
% GROUNDING DBLP:journals/ase/BanoHZT24: "qualitative research in SE has grappled with challenges like the time-intensive nature of the work, limitations to scalability due to its manual nature"
% GROUNDING DBLP:conf/chi/HeHDRH24: "The design, testing, and execution of annotation tasks took two days. In contrast, the efficiency of GPT-4 was outstanding."
The extent to which, or under what conditions, LLMs can perform these tasks \textit{effectively} remains an open research question~\cite{DBLP:conf/msr/AhmedDTP25}, and whether they should be used at all for reflexive qualitative analysis is itself contested~\cite{jowsey2025reject}. Indeed, measuring their effectiveness is practically and philosophically challenging. From an interpretivist philosophical perspective, one cannot measure the quality of analysis by comparing one (human or machine) analyst's work to another. From a realist perspective, triangulating across multiple human judges, LLM judges, and other data sources (e.g., whether a pull request is marked as having resolved an issue) improves confidence in the findings but does not prove that any one judge is valid.
% GROUNDING DBLP:conf/msr/AhmedDTP25: "we propose model-model agreement as a predictor of whether a given task is suitable for LLMs at all"
% GROUNDING jowsey2025reject: "We hold the position that analytic approaches such as reflexive thematic analysis are human research practices requiring a subjective, positioned, and reflexive researcher and therefore the use of GenAI in such approaches is not methodologically congruent."

\studytypeparagraph{Examples}
\citeauthor{Huang2023Enhancing} used multiple LLMs for joint annotation of mobile application reviews~\cite{Huang2023Enhancing}.
% GROUNDING Huang2023Enhancing: "we propose a novel approach to constructing mobile application review datasets by leveraging multiple large language models (LLMs) for joint annotation"
They used three models of comparable size with an absolute majority voting rule (i.e., a label is only accepted if it receives more than half of the total votes from the models). This approach slightly outperformed the best individual model tested. 
Meanwhile, \citeauthor{DBLP:conf/msr/AhmedDTP25} examined LLMs as annotators in SE research across five datasets, six LLMs, and ten annotation tasks~\cite{DBLP:conf/msr/AhmedDTP25}.
% GROUNDING DBLP:conf/msr/AhmedDTP25: "applying six state-of-the-art LLMs to ten annotation tasks from five datasets created by prior work"
They found that inter-model agreement strongly correlates with human-model agreement; models performed poorly in tasks where humans also frequently disagreed. They proposed to use model confidence scores to identify specific samples that could be safely delegated to LLMs, potentially reducing human annotation effort without compromising inter-rater agreement.
